%% ---- Supplementary Note [note:dango-full]: DANGO on the full trigenic dataset ----
%% Every number here is emitted by experiments/010-kuzmin-tmi/scripts/dango_full_dataset_si.py
%% (wandb pull of zhao-group/torchcell_006-kuzmin2018-tmi_dango and the three CGT replicate
%% training runs of zhao-group/torchcell_010-kuzmin-tmi_equivariant_cell_graph_transformer,
%% frozen to experiments/010-kuzmin-tmi/results/dango_full_dataset_*.csv; also writes
%% tab-dango-full-runs.tex). The figure is composed by
%% experiments/010-kuzmin-tmi/scripts/dango_full_dataset_compose_figure.py. Rerun those, not this file.
%% Dataset sizes: notes/experiments.011-kuzmin-tmi.scripts.query-comparison-006-009-010-011.md
%% (006 = 332,313; 009 = 299,146; 010 = 376,732). CGT on 009 (0.396) and 010 (0.454):
%% notes/experiments.010-kuzmin-tmi.performance-diff-010-009.md.
\clearpage
\suppnote{note:dango-full}{DANGO on the full trigenic dataset of the CGT experiment\secstatus{todo}}

The DANGO value in Fig.~\ref{fig:ggi}d, Pearson $r = 0.367$, is the mean of the three
highest of five DANGO runs trained with STRING v9.1 graphs on the trigenic dataset of the
Cell Graph Transformer (CGT) experiment; \supptab{tab:dango-full-runs} lists every run,
including nine with STRING v11.0 and v12.0 graphs, and \suppfig{fig:dango-full} shows their
training curves. Two facts qualify the comparison: the DANGO and CGT runs were trained on
two different builds of the trigenic dataset, so their validation splits are not the same
triple mutants, and every DANGO run on this dataset plateaus at $r = 0.354$ to $0.368$, well
below the $0.42$ the same implementation reaches on the Kuzmin 2018 set alone
(\suppnoteref{note:dango-repro}).

\notesec{Dataset and training} The DANGO runs use the dataset of experiment 006, the
Kuzmin 2018 trigenic interactions~\cite{kuzminSystematicAnalysisComplex2018a} with any
perturbation type plus the Kuzmin 2020 trigenic interactions restricted to deletions,
332,313 records; the CGT replicates use the dataset of experiment 010, every perturbation
type from both screens, 376,732 records. Both are split into training, validation, and test
(nominally 80/10/10) by the same seeded splitter, stratified on phenotype label and
perturbation count, but the two validation splits differ in membership, and no DANGO run on
the 010 build exists. DANGO is the implementation of \suppnoteref{note:dango-repro} (hidden
width 64, four heads) trained with the DANGO loss under the linear-to-uniform schedule
(transition at epoch 10), AdamW at learning rate $10^{-5}$, and a per-GPU batch of 64.
Twelve runs are four-GPU distributed-data-parallel (DDP) jobs on Delta A40 nodes that stop at
the 48 h queue limit after 259 to 655 epochs; two are two-GPU jobs on the IGB cluster that
complete 1,000 epochs in 59.5 h (v9.1) and 65.2 h (v12.0). The split seed is fixed and the
initialization is not, so runs of one configuration are replicates over initialization and
DDP nondeterminism. Validation Pearson $r$ is computed over the whole validation split,
synchronized across ranks; the trainer checkpoints on minimum validation MSE, and the
Pearson at that epoch is within 0.0023 of the maximum for every DANGO run.

\notesec{Results} Scored as the maximum over epochs of validation Pearson $r$, an
upward-biased order statistic reported with its epoch, the five STRING v9.1 runs reach
$0.3664 \pm 0.0004$ (mean $\pm$ SEM; range 0.3652 to 0.3676), the four v11.0 runs
$0.3557 \pm 0.0012$, and the five v12.0 runs $0.3581 \pm 0.0010$
(\suppfig{fig:dango-full}b). The three runs behind Fig.~\ref{fig:ggi}d are the v9.1 runs
\texttt{014mprap} (0.3676 at epoch 98), \texttt{x3savllr} (0.3671 at epoch 109), and
\texttt{6q08iign} (0.3664 at epoch 121); the two v9.1 runs launched later reach 0.3658 and
0.3652, so the mean over all five sits 0.0006 below the reported 0.3670. The three CGT
replicates reach $0.4537 \pm 0.0043$ at epochs 24 to 25 of 49 to 64.

\notesec{Convergence by STRING release} Every DANGO run reaches $r \ge 0.35$ within 16 to 23
epochs with v9.1, 27 to 30 with v11.0, and 25 to 30 with v12.0
(\suppfig{fig:dango-full}a,\,c), so the release delays the rise by at most ten epochs; what
it changes is the plateau. With v9.1 the maximum comes at epoch 98 to 140 and the curve then
declines (the 1,000-epoch run ends at 0.297); with v12.0 the four-GPU runs creep from 0.354
to 0.355 at epoch 100 to their maxima at epochs 336 to 608, a gain of 0.002 to 0.006, while
the two-GPU v12.0 run peaks at epoch 60 and declines. The earlier observation that DANGO
converges more slowly on v12.0 than on v9.1 holds only in this weak form, a later and lower
maximum, not a slower initial rise.

\notesec{Comparison with the Kuzmin 2018 replication} On the Kuzmin 2018 set alone (91,111
records) the same implementation reaches validation $r = 0.420$ to $0.424$ across STRING
releases (\supptab{tab:dango-string-versions}); on the 332,313-record union it reaches
$0.354$ to $0.368$. Dataset size, the split, and the addition of the Kuzmin 2020 screen
change at once, and the drop is not attributed to any one of them here. Hypothesis
(untested): the Kuzmin 2020 triples, screened around duplicated gene pairs, are the harder
part of the union, given that the CGT scores $0.396$ on the deletion-only 009 build and
$0.454$ on the 010 build. Whether DANGO trained on the 010 build would close any of the gap
to the CGT is likewise not measured; the like-for-like baseline requires that run.

\begin{figure*}[t]
\tcfig{figures/FigS-dango-full-dataset.pdf}{DANGO on the full trigenic dataset.
(a) Validation Pearson $r$ against epoch for every DANGO run on the 006 dataset, colored by
STRING release, and the three CGT replicates of Fig.~\ref{fig:ggi}d on the 010 dataset; the
epoch axis is logarithmic. (b) Maximum over epochs of validation Pearson $r$ per run; bar
and whisker are the mean and SEM over runs; filled markers are the runs averaged in
Fig.~\ref{fig:ggi}d; the axis is zoomed to 0.30 to 0.50. (c) For each DANGO run, the first
epoch at which validation Pearson reaches 0.35 (filled) and the epoch of its maximum (open);
circles are four-GPU jobs, squares two-GPU jobs. (d)--(f) Reserved: predicted against
measured interaction, absolute error against $|\tau|$, and per-channel meta-embedding
weights need a trained checkpoint and inference on the validation split; each box states the
checkpoint and script that fill it.}{fig:dango-full}
\end{figure*}

%% SOURCE: experiments/010-kuzmin-tmi/scripts/dango_full_dataset_si.py -- AUTO-GENERATED from wandb zhao-group/torchcell_006-kuzmin2018-tmi_dango and zhao-group/torchcell_010-kuzmin-tmi_equivariant_cell_graph_transformer; do not hand-edit.
\begin{table}[t]
\centering
\footnotesize
\caption{Every DANGO run on the full trigenic dataset, and the three CGT replicates of
Fig.~\ref{fig:ggi}d. Best is the maximum over epochs of validation Pearson $r$ (an
upward-biased order statistic; the epoch it occurs at follows in parentheses). At min
MSE is the validation Pearson $r$ at the epoch of minimum validation MSE, the checkpoint
the trainer keeps; it is within 0.0023 of Best for every DANGO run.
$r \ge 0.35$ is the first epoch at which the run reaches that
validation Pearson. Epochs is the number of validation epochs logged; a 4-GPU job on a
Delta A40 node stops at the 48 h queue limit, and the 2-GPU jobs on the IGB cluster ran
to 1,000 epochs. Every DANGO run used hidden width 64, four heads, the linear-to-uniform
loss schedule (transition at epoch 10), AdamW at learning rate $10^{-5}$, and a per-GPU
batch of 64, on the experiment-006 build; the CGT replicates trained on the
experiment-010 build. Runs marked $\bullet$ are the three DANGO values averaged in
Fig.~\ref{fig:ggi}d.}
\label{tab:dango-full-runs}
\begin{tabular}{@{}l l l r r r l l r l@{}}
\toprule
\textbf{Model} & \textbf{STRING} & \textbf{Run} & \textbf{GPUs} & \textbf{Epochs} & \textbf{Wall (h)} & \textbf{Best $r$ (epoch)} & \textbf{At min MSE} & $\boldsymbol{r \ge 0.35}$ & \\
\midrule
DANGO & v9.1 & \texttt{6q08iign} & 4 & 448 & 47.9 & 0.3664 (121) & 0.366 (110) & 23 & $\bullet$\\
DANGO & v9.1 & \texttt{x3savllr} & 4 & 446 & 47.9 & 0.3671 (109) & 0.367 (110) & 22 & $\bullet$\\
DANGO & v9.1 & \texttt{014mprap} & 2 & 1000 & 59.5 & 0.3676 (98) & 0.367 (81) & 16 & $\bullet$\\
DANGO & v9.1 & \texttt{cbi441o8} & 4 & 650 & 47.9 & 0.3658 (109) & 0.366 (112) & 21 & \\
DANGO & v9.1 & \texttt{p6urax0a} & 4 & 655 & 47.9 & 0.3652 (140) & 0.364 (129) & 20 & \\
DANGO & v11.0 & \texttt{ok8k5e6z} & 4 & 262 & 47.8 & 0.3536 (78) & 0.354 (67) & 29 & \\
DANGO & v11.0 & \texttt{mj7paqg8} & 4 & 259 & 47.9 & 0.3544 (76) & 0.354 (67) & 27 & \\
DANGO & v11.0 & \texttt{jo24cfpf} & 4 & 311 & 47.9 & 0.3555 (82) & 0.355 (90) & 30 & \\
DANGO & v11.0 & \texttt{3f9sh91s} & 4 & 305 & 47.8 & 0.3592 (304) & 0.359 (304) & 27 & \\
DANGO & v12.0 & \texttt{vqnrz2le} & 4 & 398 & 47.9 & 0.3582 (336) & 0.357 (326) & 25 & \\
DANGO & v12.0 & \texttt{5akd0qwt} & 4 & 446 & 47.9 & 0.3602 (432) & 0.359 (398) & 25 & \\
DANGO & v12.0 & \texttt{9jpfy547} & 2 & 1000 & 65.2 & 0.3545 (60) & 0.354 (53) & 28 & \\
DANGO & v12.0 & \texttt{xxtxdu3y} & 4 & 634 & 47.9 & 0.3597 (450) & 0.359 (433) & 30 & \\
DANGO & v12.0 & \texttt{jww4tblb} & 4 & 648 & 47.9 & 0.3577 (608) & 0.355 (90) & 28 & \\
CGT & v12.0 & \texttt{lzs9pcj3} & 4 & 64 & 62.3 & 0.4520 (24) & 0.449 (21) & 2 & \\
CGT & v12.0 & \texttt{yv4r30bi} & 4 & 64 & 62.6 & 0.4472 (25) & 0.444 (20) & 2 & \\
CGT & v12.0 & \texttt{c7671wgj} & 4 & 49 & 47.8 & 0.4619 (24) & 0.462 (24) & 2 & \\
\bottomrule
\end{tabular}
\end{table}

